% =====================================================================
% Sección: Marco teórico
% Responsable: Parte 6 — John Jairo Páez Albino
% =====================================================================
\section{Marco teórico}
\label{sec:marco}

Esta sección reúne los conceptos en los que se apoya el resto del trabajo: qué es una
red móvil \emph{ad hoc}, cómo se organiza en clústeres, cómo se modela el movimiento de
los clústeres y de los nodos, cómo encuentra rutas AODV, con qué métricas se mide el
desempeño, cómo se simula todo esto en ns-3 y cómo se analizan los resultados.

\subsection{Redes móviles \emph{ad hoc} (MANET)}
\label{sec:marco-manet}

Una red inalámbrica \emph{ad hoc} se forma sin infraestructura fija: no hay puntos de
acceso ni estaciones base, y cada nodo funciona a la vez como terminal y como enrutador,
reenviando los paquetes de otros nodos. Cuando el origen y el destino no están uno
dentro del alcance de radio del otro, el mensaje viaja en varios saltos, de vecino en
vecino. Una MANET (\emph{Mobile Ad hoc Network}) es una red \emph{ad hoc} cuyos nodos se
mueven, de modo que la topología cambia constantemente: los enlaces aparecen y
desaparecen según la distancia entre los nodos \cite{sarkar2013adhoc}.

Estas redes se despliegan cuando hace falta un servicio y desaparecen cuando este
termina, por lo que se usan en operaciones de rescate, redes militares tácticas, redes
vehiculares y lugares donde la infraestructura no existe o dejó de funcionar. Sus
desafíos principales son tres: el ancho de banda del canal inalámbrico, que comparten
todos los nodos vecinos; la energía de los dispositivos, y la inestabilidad de las rutas
causada por la movilidad \cite{sarkar2013adhoc}. Este trabajo se centra en el último.

\subsection{Arquitecturas jerárquicas y \emph{clustering}}
\label{sec:marco-jerarquia}

En una red plana todos los nodos están al mismo nivel y participan por igual en el
enrutamiento. Es sencilla, pero, a medida que crece el número de nodos, aumenta el
tráfico de control, las rutas se alargan y se vuelven inestables, y la red deja de ser
escalable. El \emph{clustering} divide la red en subconjuntos de nodos, los clústeres,
para hacerla más escalable \cite{barbeau2007principles}. Cada clúster elige una cabeza
de clúster (\emph{cluster head}, CH) que conoce a sus miembros y sabe cómo alcanzarlos.
Los miembros envían su tráfico a su CH, los CH reenvían paquetes en nombre de sus
miembros y, entre ellos, ejecutan el protocolo de enrutamiento de la red
\cite{barbeau2007principles}. Los nodos que conectan clústeres vecinos se llaman
pasarelas (\emph{gateways}) \cite{sarkar2013adhoc}. Como los CH pueden agruparse a su vez
en clústeres de un nivel superior, se obtiene una jerarquía de varios niveles; en este
trabajo los CH se comunican por un canal propio, el \emph{backbone}, y los clústeres se
agrupan en dos superclústeres (sección~\ref{sec:pregunta}).

Los clústeres pueden ser físicos o lógicos. Un clúster físico es un grupo que se
desplaza en bloque porque comparte un medio, como los pasajeros de un bus o un convoy de
vehículos. Un clúster lógico agrupa nodos por un interés común, como los integrantes de
un grupo de trabajo. En ambos casos hay dos movimientos distintos: el del clúster como
un todo y el de cada nodo dentro de él, y el modelo de movilidad debe representarlos por
separado.

\subsection{Antecedente: escalabilidad de redes \emph{ad hoc} jerárquicas}
\label{sec:marco-ospina}

Ospina y Ortiz \cite{ospina2017growth} estudiaron qué condiciones permiten que una red
\emph{ad hoc} jerárquica escale según los recursos disponibles en cada capa. Modelaron
la disponibilidad de recursos con una distribución geométrica truncada, definieron un
factor de crecimiento $j$ que relaciona los recursos de dos capas sucesivas y lo
estimaron con simulaciones en ns-3 de arquitecturas de dos y tres capas. Encontraron que
el crecimiento de los recursos entre capas sigue un modelo lineal y que el valor
estimado de $j$ está relacionado con la razón áurea. Su trabajo es un ejemplo de la
misma forma de indagación que sigue este taller: una pregunta sobre una red jerárquica,
un modelo estocástico y la simulación en ns-3 como experimento.

\subsection{Modelos de movilidad}
\label{sec:marco-movilidad}

Un modelo de movilidad describe cómo cambian en el tiempo la posición, la velocidad y
la dirección de los nodos. Los resultados de una simulación de MANET dependen mucho del
modelo elegido, así que este debe parecerse al comportamiento que se quiere estudiar
\cite{camp2002survey}. Los modelos se dividen en modelos de entidad, en los que cada
nodo se mueve de forma independiente, y modelos de grupo, en los que el movimiento de
un nodo depende del de su grupo \cite{camp2002survey}.

\paragraph{\emph{Random Walk} (caminata aleatoria).} Cada nodo escoge una dirección
uniforme en $[0, 2\pi)$ y una velocidad uniforme en $[v_{\min}, v_{\max}]$, se mueve así
durante un tiempo o una distancia fijos y vuelve a escoger; al llegar al borde del área,
rebota. El modelo no tiene memoria: el siguiente movimiento no depende del anterior, lo
que produce cambios bruscos de dirección \cite{camp2002survey}.

\paragraph{\emph{Random Waypoint} (punto de paso aleatorio).} Cada nodo escoge un
destino uniforme dentro del área y una velocidad uniforme en $[v_{\min}, v_{\max}]$,
viaja en línea recta hasta el destino, se detiene durante un tiempo de pausa y repite
el proceso. Es el modelo más usado para evaluar protocolos de MANET desde el trabajo de
Broch et al. \cite{broch1998performance}. Tiene dos defectos conocidos. El primero es
que los nodos tienden a concentrarse en el centro del área \cite{camp2002survey}. El
segundo es que, si $v_{\min} = 0$, la velocidad promedio decae con el tiempo, porque los
nodos que escogen velocidades muy bajas pasan mucho tiempo en tránsito; por eso Yoon et
al. \cite{yoon2003harmful} recomiendan una velocidad mínima positiva. Además, en lugar
de descartar el periodo inicial de la simulación, se puede arrancar desde la
distribución estacionaria del modelo \cite{navidi2004stationary}, que ns-3 ofrece como
\texttt{SteadyStateRandomWaypointMobilityModel} \cite{ns3mobility}. Las pausas también
importan: mientras un nodo está detenido no se mueve, así que la velocidad media
efectiva es menor que la velocidad configurada.

\paragraph{Gauss-Markov.} La velocidad y la dirección de cada nodo se actualizan en
intervalos fijos y conservan memoria del valor anterior. Para la velocidad,
\[
s_n = \alpha\, s_{n-1} + (1 - \alpha)\, \bar{s} + \sqrt{1 - \alpha^2}\; x_{n-1},
\]
donde $\bar{s}$ es la velocidad media, $x_{n-1}$ es una variable aleatoria normal y
$\alpha \in [0, 1]$ controla la aleatoriedad: con $\alpha = 0$ el movimiento es
completamente aleatorio y con $\alpha = 1$ es lineal. La dirección se actualiza con una
ecuación análoga \cite{camp2002survey,liang1999predictive}. El modelo evita las paradas
y los giros bruscos de los dos anteriores.

\paragraph{Movilidad de grupo por punto de referencia (RPGM).} El modelo
\emph{Reference Point Group Mobility}, propuesto por Hong et al. \cite{hong1999group},
asigna a cada grupo un centro lógico, o punto de referencia, que se mueve según un
modelo de entidad, por ejemplo \emph{Random Waypoint}. La posición de cada nodo es la
del punto de referencia más un desplazamiento aleatorio dentro de un radio. Así el grupo
se desplaza en bloque y, al mismo tiempo, cada nodo se mueve alrededor de su lugar en el
grupo.

\paragraph{Movilidad jerárquica.} La idea de RPGM se puede aplicar en varios niveles.
Si $p_S(t)$ es la posición del centro de un superclúster, $p_C(t)$ la posición del
centro de un clúster relativa a su superclúster y $p_N(t)$ la posición de un nodo
relativa a su clúster, la posición absoluta del nodo es
\[
p(t) = p_S(t) + p_C(t) + p_N(t).
\]
Cada término puede seguir su propio modelo y su propia velocidad. Por eso este enfoque
representa a la vez el movimiento del clúster y el del nodo, y permite activarlos por
separado, que es justamente lo que comparan los escenarios E1 y E2. ns-3 implementa esta
composición con \texttt{HierarchicalMobilityModel}, en la que la posición final es la
suma vectorial de la posición del modelo padre y la del modelo hijo, y ofrece
\texttt{GroupMobilityHelper} para configurar grupos completos según RPGM
\cite{ns3mobility}.

\subsection{Enrutamiento: el protocolo AODV}
\label{sec:marco-aodv}

Como la topología cambia, los protocolos de enrutamiento de las redes cableadas no son
adecuados y se usan protocolos diseñados para redes \emph{ad hoc}. Los proactivos
mantienen en todo momento rutas hacia todos los destinos mediante mensajes de control
periódicos, como OLSR \cite{clausen2003olsr}; los reactivos buscan una ruta solo cuando
hay datos que enviar, y los híbridos combinan ambos enfoques
\cite{barbeau2007principles}. Los proactivos tienen las rutas listas de inmediato, pero
gastan ancho de banda aunque no haya tráfico; los reactivos ahorran ese tráfico de
control, pero agregan una demora mientras descubren la ruta.

AODV (\emph{Ad hoc On-Demand Distance Vector}) es un protocolo reactivo
\cite{rfc3561aodv}. Cuando un nodo necesita una ruta, difunde una solicitud de ruta
(RREQ) que los vecinos retransmiten. El destino, o un nodo que conoce una ruta
suficientemente reciente hacia él, contesta con una respuesta de ruta (RREP) que regresa
por el camino inverso y deja instalada la ruta en cada salto. Si un enlace activo se
rompe, el nodo que lo detecta envía un mensaje de error (RERR) a los nodos que usaban esa
ruta, y el origen debe descubrir una nueva. Los números de secuencia de destino evitan
los ciclos y distinguen las rutas recientes de las viejas \cite{rfc3561aodv}. Por eso la
movilidad afecta directamente a AODV: cada enlace roto se traduce en paquetes perdidos o
retrasados mientras se encuentra una ruta nueva, y un enlace roto en el \emph{backbone}
afecta a todos los flujos que lo atraviesan.

\subsection{Métricas de desempeño}
\label{sec:marco-metricas}

El desempeño de la red se mide con cuatro métricas por flujo, agregadas sobre todos los
flujos de datos de una corrida (las fórmulas exactas están en la
sección~\ref{sec:pregunta}):

\begin{itemize}
    \item \textbf{Tasa de entrega de paquetes (PDR, \emph{Packet Delivery Ratio}).} Es la
    proporción de paquetes de datos que llegan a su destino respecto del total enviado
    por las aplicaciones \cite{broch1998performance}. Mide la confiabilidad de la red.
    \item \textbf{Retardo extremo a extremo.} Es el tiempo promedio entre el envío y la
    recepción de los paquetes que sí llegan. Incluye la espera en las colas, las
    retransmisiones de la capa MAC y el tiempo de descubrimiento de rutas.
    \item \textbf{\emph{Throughput}.} Es la cantidad de bits recibidos por unidad de
    tiempo durante el periodo de tráfico. Con una carga fija, baja cuando se pierden
    paquetes.
    \item \textbf{\emph{Jitter}.} Es la variación del retardo entre paquetes consecutivos
    de un mismo flujo. Un \emph{jitter} alto indica que las rutas cambian o que las colas
    crecen y se vacían de forma irregular.
\end{itemize}

En ns-3 las cuatro se calculan con el módulo \texttt{flow-monitor}, que registra para
cada flujo, entre otros contadores, los paquetes y bytes transmitidos y recibidos
(\texttt{txPackets}, \texttt{rxPackets}, \texttt{rxBytes}), la suma de los retardos
extremo a extremo de los paquetes recibidos (\texttt{delaySum}) y la suma de las
variaciones del retardo (\texttt{jitterSum}) \cite{ns3flowmonitor}.

\subsection{Simulación de eventos discretos y ns-3}
\label{sec:marco-ns3}

Una simulación de eventos discretos representa un sistema cuyo estado cambia solo en
instantes puntuales, los eventos, como el envío o la llegada de un paquete. El simulador
mantiene una lista de eventos ordenada por tiempo, toma el siguiente, actualiza el
estado y programa los eventos que este genera, sin recorrer los instantes intermedios
\cite{banks2010discrete}. Un estudio de simulación sigue estas etapas: formulación del
problema y definición de objetivos, construcción del modelo conceptual, recolección de
datos, programación, verificación y validación del modelo, diseño de experimentos,
corridas, análisis de resultados y documentación \cite{banks2010discrete}.

ns-3 es un simulador de redes basado en eventos discretos, libre y de código abierto,
escrito en C++ y orientado a la investigación y la docencia \cite{zarate2021wireless}.
Un programa de ns-3 crea nodos (\texttt{Node}), les instala dispositivos de red
(\texttt{NetDevice}) conectados a un canal (\texttt{Channel}), la pila de protocolos de
Internet con el protocolo de enrutamiento escogido, un modelo de movilidad y
aplicaciones que generan tráfico; las clases auxiliares (\emph{helpers}) simplifican
cada uno de estos pasos \cite{zarate2021wireless}. Los módulos que usa este trabajo son
\cite{ns3doc}:

{\sloppy
\begin{itemize}
    \item \texttt{wifi}: capas física y MAC de IEEE 802.11, con dispositivos en modo
    \emph{ad hoc} y canales independientes (\texttt{YansWifiChannel}).
    \item \texttt{aodv}: el protocolo AODV descrito en la sección~\ref{sec:marco-aodv}.
    \item \texttt{mobility}: los modelos de movilidad, entre ellos
    \texttt{RandomWaypointMobilityModel}, \texttt{RandomWalk2dMobilityModel} y
    \texttt{HierarchicalMobilityModel}.
    \item \texttt{internet} y \texttt{applications}: la pila TCP/IP y las aplicaciones
    de tráfico, como \texttt{OnOffApplication}.
    \item \texttt{flow-monitor}: las métricas por flujo de la
    sección~\ref{sec:marco-metricas}.
\end{itemize}\par}

El alcance de cada canal se fija con un modelo de disco
(\texttt{RangePropagationLossModel}): dos nodos tienen enlace si y solo si su distancia
no supera el alcance. Es la misma idealización que el grafo de disco unitario con que se
analizan teóricamente las redes \emph{ad hoc} \cite{barbeau2007principles}.

Las variables aleatorias de ns-3 provienen del generador MRG32k3a de L'Ecuyer
\cite{lecuyer1999good}, en una implementación que ofrece una gran cantidad de flujos y
subflujos independientes \cite{lecuyer2002object}. La semilla y el número de corrida
(\texttt{RngRun}) se fijan con la clase \texttt{RngSeedManager}; para obtener réplicas
independientes de un mismo experimento, el manual de ns-3 recomienda conservar la
semilla y cambiar el número de corrida, que selecciona otro subflujo del generador
\cite{ns3random}.

\subsection{Análisis estadístico de los resultados}
\label{sec:marco-estadistica}

El resultado de una corrida es una variable aleatoria, así que una sola corrida no basta
para responder la pregunta. Se hacen $n$ réplicas independientes de cada configuración y
se estima la media de cada métrica con un intervalo de confianza
\cite{banks2010discrete}:
\[
\bar{x} \pm t_{1-\alpha/2,\, n-1}\, \frac{s}{\sqrt{n}},
\]
donde $\bar{x}$ es la media muestral, $s$ la desviación estándar muestral y
$t_{1-\alpha/2,\, n-1}$ el cuantil de la distribución $t$ de Student con $n - 1$ grados
de libertad.

Para comparar dos escenarios, como E1 y E2, se usa la prueba $t$ de Welch
\cite{welch1947generalization}, que no supone que las dos poblaciones tengan la misma
varianza. Su estadístico y sus grados de libertad aproximados son
\[
t = \frac{\bar{x}_2 - \bar{x}_1}{\sqrt{\dfrac{s_1^2}{n_1} + \dfrac{s_2^2}{n_2}}},
\qquad
\nu \approx
\frac{\left(\dfrac{s_1^2}{n_1} + \dfrac{s_2^2}{n_2}\right)^2}
{\dfrac{(s_1^2/n_1)^2}{n_1 - 1} + \dfrac{(s_2^2/n_2)^2}{n_2 - 1}}.
\]
La prueba unilateral rechaza la hipótesis nula cuando el valor $p$ es menor que el nivel
de significancia $\alpha$; en este trabajo, $\alpha = 0.05$
(sección~\ref{sec:analisis-est}).
